\chapter{Gravitational Wave Signatures from a First-Order Phase Transition}
\label{ch:gw-signatures}

A first-order cosmological phase transition produces a stochastic
background of gravitational waves, whose spectrum encodes information
about the transition that is inaccessible to any other observational
channel. The parameters $\alpha$, $\tilde\beta$, and $v_w$ introduced
in the previous chapter determine both the amplitude and the peak
frequency of this spectrum. For an electroweak-scale transition, the
predicted signal falls precisely in the millihertz band targeted by
the LISA space-based detector, making first-order phase transitions
one of the primary science cases for next-generation gravitational
wave astronomy.

\section{Production Mechanisms}

There are three main sources of gravitational waves in a first-order
phase transition.

The first is the collision of bubble walls. As bubbles of the true
vacuum expand, their walls carry kinetic energy proportional to the
wall surface area. When two walls collide, this energy is converted
into gravitational radiation. In the envelope approximation, only the
uncollided portions of the walls contribute, and the signal is
concentrated at frequencies $f \sim \tilde\beta$.

The second and dominant mechanism in the non-runaway regime is the
generation of sound waves in the primordial plasma. As the bubble wall
expands, it pushes the plasma ahead of it, creating acoustic
perturbations that persist long after the bubble walls have collided
and dissipated. Since the energy in sound waves grows with the volume
of the bubble rather than its surface area, this mechanism dominates
over bubble wall collisions when $v_w < 1$.

The third mechanism is magnetohydrodynamic turbulence, generated when
the acoustic motion in the plasma becomes turbulent. Its contribution
is generally subdominant, at the ten percent level or less, and is
also the least well understood theoretically. We focus on the first
two mechanisms here.

\section{The Linearized Einstein Equations and the Quadrupole Formula}

We work in the weak-field limit, treating the metric as a small
perturbation about flat spacetime,
\begin{equation}
	g_{\mu\nu} = \eta_{\mu\nu} + \tilde h_{\mu\nu},
	\qquad |\tilde h_{\mu\nu}| \ll 1.
\end{equation}
Neglecting the expansion of the universe is justified because the
phase transition completes in a time $\tilde\beta^{-1} \ll H^{-1}$
much shorter than a Hubble time. Defining the trace-reversed
perturbation
\begin{equation}
	h_{\mu\nu} = \tilde h_{\mu\nu}
	- \frac{1}{2}\eta_{\mu\nu}\tilde h_{\rho\sigma}\eta^{\rho\sigma},
\end{equation}
the linearized Einstein equations in Lorenz gauge
$\partial_\mu h^{\mu\nu} = 0$ become
\begin{equation}
	\partial_\alpha\partial^\alpha h_{\mu\nu}
	= -16\pi G\,T_{\mu\nu}.
	\label{eq:linear-gw}
\end{equation}
This is a wave equation sourced by the stress-energy of the plasma and
the scalar field. It is solved by the retarded Green's function,
\begin{equation}
	G(x^\mu - y^\mu) = \frac{1}{4\pi}
	\frac{\delta(x^0 - y^0 - |\mathbf{x}-\mathbf{y}|)}
	{|\mathbf{x}-\mathbf{y}|},
\end{equation}
giving the retarded solution
\begin{equation}
	h_{\mu\nu}(t,\mathbf{x}) = -4G\int d^3y\,
	\frac{T_{\mu\nu}(t-|\mathbf{x}-\mathbf{y}|,\mathbf{y})}
	{|\mathbf{x}-\mathbf{y}|}.
\end{equation}
In the far-field limit $|\mathbf{x}-\mathbf{y}| \approx r$, using
stress-energy conservation $\partial^\mu T_{\mu\nu} = 0$ and
defining the quadrupole moment tensor
\begin{equation}
	I_{ij}(t) = \int d^3y\,T_{00}(t,\mathbf{y})\,y^i y^j,
	\label{eq:quadrupole}
\end{equation}
one can show that the gravitational wave metric perturbation is
\begin{equation}
	\boxed{h_{ij}(t,\mathbf{x})
		= -\frac{2G}{r}\frac{d^2 I_{ij}}{dt^2}(t-r).}
	\label{eq:hij}
\end{equation}
This is the gravitational analogue of the electromagnetic result that
a time-varying dipole moment sources radiation: here it is a
time-varying quadrupole moment that sources gravitational waves, which
is why spherically symmetric processes such as a star collapsing
symmetrically produce no gravitational radiation. The power radiated
is given by the quadrupole formula,
\begin{equation}
	\boxed{\frac{dE_{GW}}{dt}
		= \frac{G}{5}\frac{d^3 I_{ij}^T}{dt^3}
		\frac{d^3 I_{ij}^T}{dt^3},}
	\label{eq:quad_formula}
\end{equation}
where $I_{ij}^T$ is the traceless part of $I_{ij}$. The bubble
collisions and sound waves both produce anisotropic stress-energy that
sources a non-vanishing quadrupole moment, and hence gravitational
radiation.

\section{Order of Magnitude Estimate}

To estimate the gravitational wave signal, consider a region of size
$\sim v_w\tilde\beta^{-1}$ containing a small number of colliding
bubbles. The characteristic timescale is $\tilde\beta^{-1}$ and the
length scale is $v_w\tilde\beta^{-1}$. The kinetic energy in the
region is set by the fraction of the latent heat converted into bulk
motion,
\begin{equation}
	E_{\text{kin}} \sim \Delta\rho\,(v_w\tilde\beta^{-1})^3
	= \alpha\,\rho_{\text{rad}}\,(v_w\tilde\beta^{-1})^3.
\end{equation}
The third time derivative of the quadrupole scales as
\begin{equation}
	\frac{d^3 I_{ij}^T}{dt^3}
	\sim \frac{E_{\text{kin}}}{\tilde\beta^{-1}}
	= \alpha\,\rho_{\text{rad}}\,v_w^3\,\tilde\beta^{-2},
\end{equation}
and substituting into \eqref{eq:quad_formula},
\begin{equation}
	\frac{dE_{GW}}{dt}
	\sim G\,\alpha^2\,\rho_{\text{rad}}^2\,
	v_w^6\,\tilde\beta^{-4}.
\end{equation}
Using $G \sim H^2(T_n)/\rho_{\text{rad}}(T_n)$ from the Friedmann
equation and converting to the fractional energy density in
gravitational waves $\mathcal{S}_{GW} = \rho_{GW}/\rho_{\text{tot}}$
with $\rho_{\text{tot}} = (1+\alpha)\rho_{\text{rad}}$,
\begin{equation}
	\mathcal{S}_{GW} \sim k^2\,v_w^3\,
	\frac{\alpha^2}{(1+\alpha)^2}
	\left(\frac{H(T_n)}{\tilde\beta}\right)^2,
	\label{eq:Sgw}
\end{equation}
where $k$ is an efficiency factor, of order unity for runaway bubbles
and $\sim \alpha/(1+\alpha)$ for non-runaway cases, encoding the
fraction of latent heat that goes into bulk kinetic energy rather than
reheating the plasma. The characteristic frequency of the signal at
production is $f_{GW} \sim \tilde\beta$, set by the inverse duration
of the transition.

Several features of \eqref{eq:Sgw} deserve comment. The signal is
enhanced by large $\alpha$, a strongly supercooled transition, and
suppressed by large $\tilde\beta/H$, a short transition. The
dependence on $v_w$ reflects the fact that faster walls convert more
latent heat into kinetic energy before the walls collide. The
prefactor $(H/\tilde\beta)^2 \ll 1$ generically suppresses the signal,
which is why only the strongest transitions produce a detectable
signal.

\section{The Spectrum Today}

The gravitational wave spectrum observed today is obtained by
redshifting from the time of the transition at temperature $T_n$ to
the present. Using entropy conservation,
\begin{equation}
	g_{*s,n}\,T_n^3\,a_n^3 = g_{*s,0}\,T_0^3\,a_0^3,
\end{equation}
together with the Friedmann equation
$H^2 = \rho_{\text{tot}}/3M_{\text{pl}}^2$ with
$M_{\text{pl}} = 2.435\times 10^{18}\,\text{GeV}$, and the
present-day values $H_0 = h\cdot 2.13\times 10^{-42}\,\text{GeV}$,
$T_0 = 2.43\times 10^{-13}\,\text{GeV}$, and $g_{*s,0} = 3.97$,
the present-day fractional gravitational wave energy density is
\begin{equation}
	\Omega_{GW,0}\,h^2 \sim 1.89\times 10^{-5}\,k^2\,v_w^3\,
	\frac{\alpha^2}{(1+\alpha)^2}
	\left(\frac{H(T_n)}{\tilde\beta}\right)^2
	\left(\frac{g_{*s,n}}{100}\right)^{-1/3},
	\label{eq:Omega_gw}
\end{equation}
and the peak frequency today is
\begin{equation}
	f_{GW,0} \sim 1.7\times 10^{-5}\,\text{Hz}\,
	\left(\frac{\tilde\beta}{H(T_n)}\right)
	\frac{T_n}{100\,\text{GeV}}
	\left(\frac{g_{*s,n}}{100}\right)^{1/6}.
	\label{eq:f_gw}
\end{equation}
For an electroweak-scale transition with $T_n \sim 100\,\text{GeV}$
and $\tilde\beta/H \sim 10^2$, equation \eqref{eq:f_gw} gives a peak
frequency of order $10^{-3}\,\text{Hz}$, which falls squarely in the
sensitivity band of the LISA space interferometer, planned for launch
in the 2030s. The amplitude \eqref{eq:Omega_gw} for a moderately
strong transition with $\alpha \sim 0.1$ and $k \sim 0.1$ is of order
$\Omega_{GW,0}h^2 \sim 10^{-12}$, at the edge of LISA's sensitivity.
Stronger transitions with larger $\alpha$ or smaller $\tilde\beta/H$
produce correspondingly larger signals.

These results illustrate the power of gravitational wave astronomy as
a probe of the early universe. The spectrum \eqref{eq:Omega_gw} is
determined entirely by the thermodynamic parameters of the phase
transition, which in turn are calculable from the finite-temperature
effective potential developed in the preceding chapters. A detection
by LISA would not only confirm the occurrence of a first-order
electroweak phase transition but would also constrain the parameters
of the underlying particle physics model, providing information about
physics at the electroweak scale that is complementary to and in some
respects beyond what is accessible at the LHC.
